\section{Thesis Structure}
\label{sec:thesis_structure}

The remainder of this report is organised into six chapters as follows.

\textbf{Chapter 2 — Theoretical Background} reviews the foundational concepts and prior work relevant to this project.
This includes the theoretical basis of sequential recommendation models, graph-based collaborative filtering, dense retrieval with transformer encoders, approximate nearest-neighbour search, and multilingual language modelling.
Existing commercial and academic recommendation systems are also surveyed to contextualise the proposed approach.

\textbf{Chapter 3 — System Analysis} describes the design and prior implementation of the data platform.
The overall architecture, individual module designs (ingestion, visualisation, storage), and the legacy GRU-SeqDQN recommendation component are analysed.
The chapter concludes with an identification of the system's limitations that motivate the proposed improvements.

\textbf{Chapter 4 — Proposed Solutions} presents the technical contributions of this project.
The infrastructure changes are described first, followed by the improvements to each of the three data platform modules.
The dual-mode recommendation architecture — Pipeline~A (Cleora + BGE-M3) and Pipeline~B (DIF-SASRec) — is introduced and justified in detail.

\textbf{Chapter 5 — Experiments and Evaluation} reports the quantitative evaluation of the proposed system.
Component-level analyses (encoder comparison, HNSW vs.\ flat index, search latency profiling) are presented first, followed by end-to-end recommendation results on 100{,}000 users, sensitivity analyses, and a qualitative pipeline complementarity study.

\textbf{Chapter 6 — Application} demonstrates the generality of the data platform by applying it to an air quality monitoring use case.
The chapter covers the domain problem, dataset description, data ingestion and preparation, visualisation, retrieval, and a summary of results achieved in this domain.

\textbf{Chapter 7 — Conclusion} summarises the achievements of the project against the stated goals, identifies limitations, and proposes directions for future work.

